\lyseis{Β1.11}{Ο κύκλος και τα στοιχεία του}
\sol{1} Ελεύθερη κατασκευή.
\sol{2} α) $8$ \quad β) $4{,}5$ \quad γ) $5$ \quad δ) $3{,}6$.
\sol{3} Το $\mathrm{Α}$ μέσα, το $\mathrm{Β}$ πάνω, το $\mathrm{Γ}$ έξω, το $\mathrm{Δ}$ (στο κέντρο) μέσα.
\sol{4} Η μεγαλύτερη χορδή είναι η διάμετρος, $10$. Υπάρχουν άπειρες διάμετροι.
\sol{5} Ελεύθερη κατασκευή.
\sol{6} α) Σωστό. β) Λάθος: η διάμετρος. γ) Σωστό. δ) Λάθος: είναι μέσα. ε) Σωστό.
\sol{7} $5$ μ.
\sol{8} $33$ cm. Ο άξονας είναι στα $33$ cm από το έδαφος (η ακτίνα).
\sol{9} Σε κύκλο ακτίνας $8$ cm, διαμέτρου $16$ cm.
\sol{10} $18{,}3$ μ.
\sol{11} Όχι: η διάμετρος της πίτας ($26$ cm) είναι μεγαλύτερη από την πλευρά του ταψιού ($25$ cm).
\sol{12} Γιατί για κάθε χορδή το τμήμα $\mathrm{ΑΒ}\le\mathrm{ΑΟ}+\mathrm{ΟΒ}=2\rho$ (η τεθλασμένη $\mathrm{ΑΟΒ}$ είναι μεγαλύτερη ή ίση από το $\mathrm{ΑΒ}$), και ίσο μόνο αν $\mathrm{Ο}$ στο $\mathrm{ΑΒ}$.
\sol{13} Δύο: ο κύκλος με κέντρο το σημείο και ακτίνα $5$ τέμνει τον αρχικό σε δύο σημεία (αν το σημείο είναι στον κύκλο).
\sol{14} Άπειροι. Τα κέντρα βρίσκονται στη μεσοκάθετο του $\mathrm{ΑΒ}$ (θα το δούμε στο επόμενο κεφάλαιο).
\sol{15} Όχι: ένα κοινό σημείο θα είχε την ίδια απόσταση από το κέντρο, άρα οι ακτίνες θα ήταν ίσες και οι κύκλοι θα συνέπιπταν.
\sol{16} Ελεύθερη απάντηση.
\sol{17} Ελεύθερη απάντηση.
